\section*{Chapitre \ref{ch:s2:swap-spread-and-asset-swap-trades} --- Trades de swap spread et d'asset swap}

\begin{solution}{exo:s2:swap-spread-and-asset-swap-trades:1}
$0.05 \times 0.10 = 0.005$ du notionnel par an~: 50 points de base.
\end{solution}

\begin{solution}{exo:s2:swap-spread-and-asset-swap-trades:2}
$8.5 \times 12 = 102$ points de base de notionnel.
\end{solution}

\begin{solution}{exo:s2:swap-spread-and-asset-swap-trades:3}
$40 + 5 - 50 = -5$ points de base par an~: la charge l'emporte sur le carry.
\end{solution}

\begin{solution}{exo:s2:swap-spread-and-asset-swap-trades:4}
Une position à deux ans mobilise peu de bilan pour une faible duration, et la prime de crédit bancaire dans la jambe variable compte encore aux échéances courtes~;
à trente ans, le carry par unité de bilan est faible au regard de la duration couverte, et les receveurs de taux à long terme (fonds de pension, assureurs)
font baisser le taux de swap.
\end{solution}

\begin{solution}{exo:s2:swap-spread-and-asset-swap-trades:5}
Rien ne change, ni côté crédit ni côté duration, les derniers jours d'un trimestre~; ce qui change, c'est ce que les banques déclarent. Une baisse ces jours-là traduit le coût de
détenir des positions à la date de déclaration.
\end{solution}

\begin{solution}{exo:s2:swap-spread-and-asset-swap-trades:6}
Les détenteurs dont le \omterm{def:s2:swap-spread-and-asset-swap-trades:bscost}{coût de bilan} propre est inférieur au carry~: fonds disposant de repo à terme bon marché, investisseurs détenant des bons du Trésor en direct, ou tout acteur dont le capital est
mesuré sur la même base. Pour eux, le carry de 38.8 points de base par an est surtout un profit.
\end{solution}

\begin{solution}{exo:s2:swap-spread-and-asset-swap-trades:7}
Sur les dix ans, il perd 36.8 points de base par an~: le carry rapporte 21.7, mais les réévaluations font perdre 63.4, car le spread est passé de sa prime à environ
$-39$ à mesure que le \omterm{def:s2:swap-spread-and-asset-swap-trades:bscost}{coût de bilan} s'est imposé. Un trade de carry ouvert avant que son moteur ne change paie en valorisations la totalité du réajustement.
\end{solution}

\begin{solution}{exo:s2:swap-spread-and-asset-swap-trades:8}
Elle n'est pas sans risque~: la position doit être financée en repo et consomme du bilan (50 points de base par an dans la charge du chapitre), et le spread de swap
peut encore se creuser, valorisé à un DV01 de 8.5. Ce qui ressemble à un arbitrage est le prix du bilan qu'elle exige.
\end{solution}

\begin{solution}{pb:s2:swap-spread-and-asset-swap-trades:1}
\begin{enumerate}
\item Taux de swap moins rendement de l'emprunt d'État~; le spread sur le taux variable que verse une obligation swappée.
\item Une prime pour le crédit bancaire dans la jambe variable et pour la liquidité des bons du Trésor.
\item Dix ans~: 58.6 avant et 8.3 après~; trente ans~: 48.3 avant et $-21.4$ après.
\item Les frictions liées à la détention d'obligations rendent les spreads négatifs peu surprenants.
\item Le rendement exigé sur le capital immobilisé face aux actifs bruts d'un trade, exprimé en taux courant.
\item Écarts persistants de parité des taux d'intérêt couverte (CIP), maximaux en fin de trimestre, qui renvoient à la réglementation.
\item Il recule d'environ 11.5 à environ $-38.8$ points de base à mesure que le coût s'impose progressivement.
\item Les cinq derniers jours de chaque trimestre~; deux fois plus profond en fin d'année~: dates de déclaration.
\item Position longue sur l'obligation financée en repo, payeur de taux fixe~; carry, réévaluations, variable moins repo, charge.
\item $-24.3$ et $-74.3$ avant~; 32.4 et $-17.6$ après.
\item Le bruit du spread, à un DV01 de 8.5, est grand face au carry.
\item 103.3 points de base en fin d'année et 43.9 aux autres fins de trimestre.
\item \textbf{Résultat nommé.} Une fois le \omterm{def:s2:swap-spread-and-asset-swap-trades:bscost}{coût de bilan} intégré, le trade long de swap spread rapporte 32.4 points de base par an hors charge et perd 17.6
points de base par an une fois la charge de 50 points de base par an appliquée.
\item Là où le détenteur marginal, une banque, est indifférent après sa charge.
\item Un bilan moins coûteux (un changement de règles ou de détenteurs des bons du Trésor).
\item Par la volatilité du spread multipliée par le DV01 et par le bilan disponible.
\item Facturer le bilan, intégrer le changement de régime, utiliser le repo réellement disponible.
\item Le carry d'asset swap ou le spread de swap long.
\item Ce sont deux positions sur bons du Trésor financées en repo, dont le rendement dépend du financement et du collatéral.
\item Le coût de détention des bons du Trésor au bilan d'une banque.
\end{enumerate}
\end{solution}

\begin{solution}{iq:s2:swap-spread-and-asset-swap-trades:1}
Taux de swap moins rendement de l'emprunt d'État à même échéance~; une valeur négative signifie que l'obligation d'État rapporte plus que le swap, ce qu'un détenteur ne peut verrouiller
qu'en finançant l'obligation, ce qui coûte du repo et du bilan.
\end{solution}

\begin{solution}{iq:s2:swap-spread-and-asset-swap-trades:2}
Détenir l'obligation consomme du bilan que les banques facturent, alors que recevoir le taux fixe sur le swap en consomme peu~; les receveurs de taux à long terme maintiennent les taux de swap bas.
Le modèle de Jermann, avec des frictions sur la détention d'obligations, produit des spreads négatifs.
\end{solution}

\begin{solution}{iq:s2:swap-spread-and-asset-swap-trades:3}
Acheter l'obligation, payer le taux fixe sur un swap aux coupons identiques et recevoir le taux variable plus un spread~: le spread d'asset swap. Payer le taux de repo pour financer
l'obligation et le \omterm{def:s2:swap-spread-and-asset-swap-trades:bscost}{coût de bilan} de sa détention.
\end{solution}

\begin{solution}{iq:s2:swap-spread-and-asset-swap-trades:4}
Un nouvel élargissement vers des valeurs négatives (une hausse du \omterm{def:s2:swap-spread-and-asset-swap-trades:bscost}{coût de bilan}), un pic du repo, une baisse de fin de trimestre plus marquée qu'à l'accoutumée, un changement d'indice variable, et
des hausses des marges de compensation.
\end{solution}

\begin{solution}{iq:s2:swap-spread-and-asset-swap-trades:5}
Allouer le capital que chaque position exige selon le ratio contraignant, le facturer chaque jour au taux de rendement minimal exigé (hurdle rate) de la firme, et présenter le P\&L des desks après la
charge, avec des taux plus élevés pour les positions conservées au-delà des dates de déclaration.
\end{solution}

\begin{solution}{iq:s2:swap-spread-and-asset-swap-trades:6}
L'obligation rapporte son rendement et coûte le repo~; le swap reçoit le taux variable et verse le taux fixe~; la somme donne rendement moins taux de swap plus taux variable moins
repo. Si l'indice variable s'écarte du repo (comme lorsqu'un taux interbancaire est remplacé par un taux garanti), le dernier terme change et avec lui le
carry du trade.
\end{solution}
